\documentclass{article}
\usepackage[utf8]{inputenc}
\usepackage{geometry}
\usepackage{parskip}
\usepackage{hyperref}
\geometry{a4paper, margin=1in}

\title{Project Development Diary and Retrospective}
\author{Botir Khaltaev}
\date{}

\begin{document}

\maketitle

\section{Introduction}

This document presents a retrospective analysis of the development of the TensorStore project during the first term. It compares the original project plan and timeline against the work actually carried out, reflects on deviations from the plan, evaluates how identified risks materialised in practice, and outlines a structured and motivated plan for the second term. The focus is on technical decision making, benchmarking discipline, project management, and lessons learned from real system development.

\section{Comparison with the Original Plan and Timeline}

The original project plan divided the work into four phases: background research during Weeks 1 to 2, core validation during Week 3, analysis and documentation during Week 4, and advanced features during Weeks 5 to 15. This structure was intentionally designed to enforce early validation of the central hypothesis that Linux \texttt{io\_uring} could significantly outperform synchronous and multi-threaded I/O when loading large LLM checkpoints.

Phase 1 progressed largely as planned. Background research into ServerlessLLM architectures, checkpoint formats, and the Linux I/O stack was completed on schedule. During this phase, the \texttt{io\_uring} programming model and available Rust bindings were evaluated, and an initial TensorStore design was established based on this research.

Phase 2 was intended to act as a strict validation checkpoint. Early benchmarks compared synchronous file I/O, Tokio-based asynchronous I/O, and \texttt{io\_uring}-based loading. These results showed only a modest improvement of approximately five percent when using \texttt{io\_uring}. According to the original project plan, this outcome should have triggered either a significant narrowing of scope or a redefinition of the project goals. In practice, development continued beyond this checkpoint without a formal decision being documented.

Phase 3 did not occur as cleanly as planned. While some performance analysis was carried out, documentation and structured evaluation lagged behind implementation. Meaningful analysis only became possible after further refactoring, correctness fixes, and improvements to the benchmarking methodology were applied later in the term.

Phase 4 work began significantly earlier than planned and expanded beyond its original scope. During the first term, features such as buffer pooling, zero-copy loading paths, direct I/O support, comprehensive automated testing, continuous integration, and full checkpoint format conversion pipelines were implemented. While this reduced adherence to the original phase boundaries, it resulted in a substantially more robust experimental system than initially planned.

By the end of Term 1, the project had progressed much further in implementation depth than originally scheduled, but at the cost of delayed hypothesis validation and documentation.

\section{Reflection on Term 1 Execution}

Several aspects of the project execution worked well. Working with real-world LLM checkpoints exposed performance bottlenecks and correctness issues that would not have appeared in synthetic benchmarks. Investing in automated testing and continuous integration significantly reduced regression risk and increased confidence in experimental results. Repeated refactoring also led to a clearer and more maintainable system architecture.

However, there were notable shortcomings. The most significant issue was a lack of discipline at the planned validation checkpoint. Once early benchmarks failed to demonstrate the expected performance improvement, a formal reassessment of scope should have occurred. Instead, development gradually shifted toward feature-driven engineering, driven by the desire to better understand the system rather than to strictly evaluate the original hypothesis.

Another major mistake involved the initial benchmarking methodology. Early benchmarks unintentionally factored in cache warmness at multiple levels, including the Linux page cache and internal buffer reuse. As a result, these measurements did not accurately represent the cold-start scenario that is central to serverless LLM systems, where models are loaded on demand with no prior state.

This led to misleading conclusions, as cached data masked the true cost of checkpoint loading and reduced the apparent performance differences between I/O strategies. After identifying this issue, the benchmarking approach was revised to measure cold starts only. Caches were cleared between runs, buffer reuse was restricted, and each benchmark execution was treated as an isolated model load. This change significantly improved the validity of the evaluation and aligned it more closely with real-world deployment conditions.

Documentation also lagged behind code development until late in the term, increasing cognitive overhead and slowing iteration. This reinforced the importance of writing alongside implementation rather than treating documentation as a final step.

The primary lesson learned from Term 1 is that hypothesis-driven systems research requires enforced stopping points and carefully controlled evaluation. Without these constraints, development can drift toward building increasingly complex systems without sufficient analytical grounding.

\section{Risk Assessment and Mitigation}

Most of the risks identified in the original project plan materialised during Term 1.

The risk of hypothesis invalidation became evident when early benchmarks failed to show a large performance advantage for \texttt{io\_uring}. This was mitigated by reframing the project to focus on identifying the true bottlenecks in checkpoint loading rather than solely validating a single I/O mechanism.

Ecosystem maturity also proved to be a genuine risk. Limited activity and documentation in the \texttt{tokio-uring} ecosystem required conservative use of advanced features and extensive correctness testing. This risk was mitigated through modular backend design and comprehensive automated testing.

Time estimation risk was partially mitigated through incremental development and task decomposition. However, scope expansion still exceeded original estimates, highlighting the need for stricter milestone enforcement and clearer decision points in future work.

Overall, the risks identified in the project plan were appropriate and relevant, but their mitigation would have benefited from more explicit enforcement of planned evaluation and decision stages.

\section{Term 2 Plan and Forward Milestones}

The second term will prioritise consolidation, structured evaluation, and documentation over further feature expansion. Planned milestones are as follows.

During Weeks 1 to 3, the focus will be on formal performance evaluation using controlled cold-start benchmarks. This will include detailed CPU, memory, and I/O profiling to isolate bottlenecks identified during Term 1.

During Weeks 4 to 6, TensorStore will be integrated into an end-to-end inference pipeline. This will allow measurement of user-visible cold-start latency rather than relying solely on synthetic throughput metrics.

During Weeks 7 to 9, documentation and analysis will be developed in parallel with experimentation. This includes structured result interpretation, comparison against baseline approaches, and preparation of figures and tables for the final report.

During Weeks 10 to 12, a final evaluation will be conducted to assess whether the performance benefits of \texttt{io\_uring} justify its added complexity in practical serverless inference systems. This evaluation will directly inform the conclusions of the final report.

These milestones directly address the weaknesses identified during Term 1 and aim to ensure a disciplined and evidence-driven conclusion to the project.

\section{Conclusion}

This retrospective demonstrates a clear comparison between the original project plan and the actual execution of the TensorStore project during Term 1. While the project deviated from its intended hypothesis-driven timeline, it produced valuable technical insights and a robust experimental framework. The lessons learned regarding benchmarking discipline, validation checkpoints, and documentation practices provide a strong foundation for a more structured and evaluative approach in Term 2.

\end{document}
